% Original mechanism-based training problem, review edition.
\documentclass[11pt,a4paper]{article}
\usepackage[margin=23mm,headheight=14pt]{geometry}
\usepackage{fontspec}
\setmainfont{Noto Serif CJK SC}[ItalicFont={Noto Serif CJK SC},ItalicFeatures={FakeSlant=0.15}]
\setsansfont{Noto Sans CJK SC}\setmonofont{DejaVu Sans Mono}[Scale=0.81]
\XeTeXlinebreaklocale "zh"\XeTeXlinebreakskip=0pt plus 1pt
\usepackage{amsmath,amssymb,booktabs,array,enumitem,tikz,fancyhdr,hyperref,xurl}
\hypersetup{unicode=true,colorlinks=true,linkcolor=black,urlcolor=blue!50!black,pdftitle={轻型屋盖的积水与排水},pdfauthor={SZUMMA2026 Training}}
\setlength{\parindent}{2em}\setlength{\parskip}{5pt}
\setlist[enumerate]{leftmargin=2em,itemsep=11pt,topsep=7pt}
\pagestyle{fancy}\fancyhf{}\fancyhead[L]{\small 数学建模训练题\quad Problem\_004}\fancyhead[R]{\small 审核稿}\fancyfoot[C]{\thepage}
\renewcommand{\headrulewidth}{0.3pt}\newcommand{\file}[1]{\texttt{\detokenize{#1}}}
\begin{document}
\begin{center}{\LARGE\bfseries 轻型屋盖的积水与排水}\end{center}

平屋盖排水不畅时，积水会使屋面下挠，原来的低洼处可能进一步加深，排水口与屋面高点的位置也随之改变。只按未变形的屋面计算雨水荷载，未必能反映实际情况。某场馆准备改造一段轻型屋盖，需比较排水改造、梁加固和临时顶撑的作用。

取屋盖中长 $L=6$ 米、宽 $b=4$ 米的一条带，$x$ 从左端沿长度方向量起。本题假定同一 $x$ 处沿宽度方向的各点升降相同，不考虑横向弯曲。因此，可用一根长6米的等效梁描述整个条带的竖向变形，荷载和刚度均按整个4米宽条带计算。

梁的左右端点由支座固定在同一高度，不能上下移动，但梁在端部可以倾斜，支座不限制它的转动，也不施加弯矩。端部设有锚固，梁不能离开支座；支座可对梁施加向上或向下的竖向力。梁本身可以在荷载作用下弯曲。

以梁未受荷载时的水平位置为高程零线。挠度 $w(x)$ 是梁在 $x$ 处相对零线的竖向位移，本题向下取正，故变形后屋面高程为 $-w(x)$。采用小挠度线弹性梁模型：梁的转角较小，材料在弹性范围内工作，忽略剪切变形及梁沿长度方向的伸长。

抗弯刚度 $EI$ 表示梁抵抗弯曲的能力，数值越大，同样荷载下越不易弯曲。$E$ 为材料的弹性模量，$I$ 为描述截面形状对弯曲影响的截面二次矩；本题已给出两者的乘积，无须另查材料或截面尺寸。永久荷载指屋面及构件等始终存在的重量，以均布线荷载 $q_0$ 表示：长为 $\Delta x$ 的一小段承受向下的力 $q_0\Delta x$，不应再乘条带宽度。初始无水时，梁也已因永久荷载而下挠。

\begin{center}
\begin{tikzpicture}[x=1.65cm,y=1.25cm,font=\small,>=stealth]
\draw[dashed] (0,0)--(6,0);
\draw[very thick,domain=0:6,samples=61] plot (\x,{-0.08*\x*(6-\x)});
\foreach \x in {.65,1.4,2.15,3,3.85,4.6,5.35}{
  \draw[->] (\x,.6)--(\x,{-0.08*\x*(6-\x)+.07});
}
\node[above] at (3,.62) {向下的分布荷载 $q(x)$，单位 N/m};
\draw (0,0)--(-.13,-.23)--(.13,-.23)--cycle;
\draw (6,0)--(5.87,-.23)--(6.13,-.23)--cycle;
\draw (-.21,-.27)--(.21,-.27);\draw (5.79,-.27)--(6.21,-.27);
\node[left] at (0,.03) {$x=0$};\node[right] at (6,.03) {$x=L$};
\draw[<->] (2.55,0)--(2.55,{-0.08*2.55*(6-2.55)}) node[midway,left,fill=white,inner sep=1pt] {$w(x)$};
\node[anchor=east] at (6,.16) {未受荷载时的零线};
\node at (3,-1.00) {变形后的梁（竖向变形放大）};
\draw[<->] (0,-1.40)--(6,-1.40) node[midway,fill=white,inner sep=2pt] {$L=6\,\mathrm m$};
\node at (3,-1.79) {两端均不能上下移动，但允许梁的端部转动};
\end{tikzpicture}
\end{center}

为便于使用上述梁模型，给出其基本关系。在没有集中力作用的位置，若合计竖向分布荷载为 $q(x)$，向下取正，则
\[
EI\frac{\mathrm d^4w}{\mathrm dx^4}=q(x),\qquad
M(x)=-EI\frac{\mathrm d^2w}{\mathrm dx^2}.
\]
把梁在某处假想切开，截面两侧除了相互传递力，还会传递使梁弯曲的转动作用，称为弯矩 $M(x)$，单位为牛米。弯矩的正负表示相反的弯曲方向，本题均以 $|M(x)|$ 与限值比较。上述支座对应的端部条件为
\[
w(0)=w(L)=0,\qquad M(0)=M(L)=0.
\]
式中的 $q(x)$ 应包括永久荷载和积水作用。顶撑施加在一个位置的力属于集中力，须单独计入受力平衡。

\newpage
\begin{center}{\large\bfseries 积水分区与排水}\end{center}

中点有一道随屋面运动的隔坎，将屋面分成左右两区。每区内水面水平，水位分别为 $H_1$、$H_2$；水深由水面与变形后屋面的高差确定，干处水深为零。两区仅通过隔坎上的一个槽口交换水。设坎高为 $d$、槽宽为 $b_c$，坎顶高程为 $z_c=d-w(3)$。本题按下式计算过坎流量：
\[
Q_{12}=C_w b_c\left([H_1-z_c]_+^{3/2}-[H_2-z_c]_+^{3/2}\right),
\qquad [a]_+=\max(a,0).
\]
$Q_{12}>0$ 表示从左区流向右区。该式是题设采用的分区交换关系，不要求另建淹没出流修正。

两区各有一个排水口，随当地屋面运动。水须高于排水口的进水下缘（称为唇口）才能排出；唇口高出当地屋面 $e$，因此唇口高程为 $-w(x_j)+e$。两口的位置、有效面积分别为 $x_j,A_j$。下游通大气，无回流，排水量按
\[
Q_j=C_d A_j\sqrt{2g[H_j+w(x_j)-e]_+},\qquad j=1,2
\]
计算。流量公式均用国际单位制。题给进水量已包含本条带降雨及邻接屋面的汇水，不再另加降雨。忽略蒸发、渗漏及外缘溢流，停雨后排水仍继续。

\begin{center}
\begin{tikzpicture}[x=1.6cm,y=.9cm,font=\small,>=stealth]
\draw[thick] (0,0) rectangle (6,1.7);
\draw[very thick] (3,0)--(3,.68);\draw[very thick] (3,1.02)--(3,1.7);
\node at (1.5,1.25) {左区};\node at (4.5,1.25) {右区};
\draw[fill=white] (1.2,.65) circle (.06);\draw[fill=white] (4.8,.65) circle (.06);
\node[below] at (1.2,.52) {排水口1};\node[below] at (4.8,.52) {排水口2};
\draw[->] (3.7,2.08)--(3,.85);\node[right] at (3.7,2.08) {隔坎槽口};
\draw[<->] (-.36,0)--(-.36,1.7) node[midway,left] {$b=4\,\mathrm m$};
\draw[<->] (0,-.48)--(6,-.48) node[midway,fill=white,inner sep=2pt] {$L=6\,\mathrm m$};
\node at (3,-.98) {平面布置示意，隔坎位于 $x=3$ 米处};
\end{tikzpicture}
\end{center}

\begin{center}
\begin{tikzpicture}[x=1.65cm,y=2cm,font=\small,>=stealth]
\draw[dashed] (0,0)--(6,0);\node[above] at (3,.06) {零线（未加载屋面）};
\draw[very thick,domain=0:6,samples=61] plot (\x,{-0.08*\x*(6-\x)});
\draw[very thick] (3,-.72)--(3,-.30);
\draw[blue,thick] (.5762,-.25)--(2.97,-.25);\draw[blue,thick] (3.03,-.40)--(5,-.40);
\node[blue,above] at (1.7,-.25) {水面 $H_1$};\node[blue,above] at (4.5,-.40) {水面 $H_2$};
\draw[<->] (2.82,-.72)--(2.82,-.30) node[midway,left] {$d$};
\draw[dotted] (2.76,-.72)--(3.05,-.72);\draw[dotted] (2.76,-.30)--(3.05,-.30);
\draw[fill=white] (1.2,-.4608) circle (.04);\draw[fill=white] (4.8,-.4608) circle (.04);
\draw[->] (1.0,-.99)--(1.2,-.51);\node[below] at (1.0,-.99) {排水口1};
\draw[->] (5.0,-.99)--(4.8,-.51);\node[below] at (5.0,-.99) {排水口2};
\draw[->] (3.6,-1.00)--(3,-.70);\node[below] at (3.6,-1.00) {隔坎随屋面升降};
\node at (3,-1.47) {纵剖面示意：水位相对零线量取，水深从当地屋面量取};
\end{tikzpicture}
\end{center}

水面高程 $H_j$ 与当地水深是不同的量：水面可在零线以下而仍有积水，水深应从已经下挠的屋面量起。图中水位低于零线仅为示意，并非规定两区水位必须为负。

结构变形与区内水面调整远快于进排水过程。因此，本题不计算梁上下振动的过程，而是在每个时刻将梁看作已在当前积水作用下达到静力平衡，即准静态处理。忽略水的附加惯性及水平水压对等效梁的影响。各图不按比例绘制，也不代表给定时刻的计算结果，具体参数见附件。

\newpage
\begin{center}{\large\bfseries 研究任务}\end{center}

请建立数学模型，研究以下问题。
\begin{enumerate}
\item 暂时拆除隔坎、封住排水口，比较两种静水试验：一种加入有限水量后封闭，另一种与大水箱双向连通以维持给定水位。先对总水量为 $0.1、0.5、1.0\,\mathrm{m^3}$ 的封闭试验，计算原梁的变形、水面及弯矩。再以维持水位 $H=0.01\,\mathrm m$ 为例，分析改变 $EI$ 时的平衡与稳定性。两种试验的稳定性判断有什么不同？说明所用判据，并区分稳定平衡与满足使用限值。
\item 恢复隔坎和原排水口。附件给出 U、L、R 三种进水工况，总进水流量随时间完全相同，左右分配不同。分别研究从干屋面开始的90分钟过程，分析变形、积水、分区交换及排水。总蓄水量是否足以评价屋盖状态？将屋面固定在永久荷载下的初始形状，重新计算进排水，比较忽略水荷载变形对判断的影响。给出检验模型和数值结果的依据。
\item 在1500元、7500元两种预算下，分别提出永久改造方案。可选加固档位、排水口位置与有效面积见附件，同一方案须用于三种工况。整个过程向下挠度不超过 $0.12\,\mathrm m$，正负弯矩的绝对值不超过所选档位的限值，90分钟末两区总残余水量不超过 $0.18\,\mathrm{m^3}$。比较费用、使用余量与排水结果，说明选择依据。
\item 采用第三问7500元预算推荐方案，若左排水口在第25分钟完全堵塞，此后不再排水，而右口保持正常，拟在中点设置临时顶撑。顶撑只向上承压、不向下拉，运动规则如下。比较不顶撑和不同启动时刻的过程，给出一个同时适用于三种进水工况的启动安排，或说明不能满足要求的原因。须保留第三问的挠度、正负弯矩限值，顶撑反力不超过5000牛；启动前已出现的越限不能由后来恢复抵消。考察至90分钟，报告残余水量，但不以第三问残余水量限值作为应急通过条件。评价减小挠度是否足以说明恢复了排水，顶撑能否代替疏通排水口。
\end{enumerate}

顶撑启动时刻 $t_J$ 可在第30至60分钟内选择。启动前支柱不接触屋盖。令
\[
s_0=\max\{w(3,t_J^-),0.012\},\qquad
s(t)=\max\{0.012,s_0-0.002(t-t_J)\},\quad t\ge t_J,
\]
其中时间以分钟计，长度以米计。支柱顶高程为 $-s(t)$，即以每分钟2毫米向上移动，达到距零线向下12毫米处后停止。中点屋面不能向下穿过支柱；若屋面与支柱分离，则顶撑反力为零。临时顶撑不改变梁刚度、隔坎或排水口参数，考核其反力，不计其增设费用。

论文应说明模型假设、稳定性与使用限值的关系，附可复现的程序。题目参数与限值用于方案比较，不视为现行建筑规范条款。

\newpage
\begin{center}{\large\bfseries 参数与附件}\end{center}

本题附件全部为设计参数与规定工况，不是实测记录。研究按上述梁模型开展，无须另行收集监测数据。水面高程、局部水深和平均水深应加以区分。

\begin{center}
\begin{tabular}{lrl}
\toprule
项目 & 数值 & 说明及单位\\\midrule
长度、宽度 & 6、4 & m\\
原梁 $EI_0$ & $4\times10^5$ & 整条带等效抗弯刚度，$\mathrm{N\,m^2}$\\
$q_0$ & 600 & 整条带永久线荷载，N/m\\
$\rho$、$g$ & 1000、9.81 & $\mathrm{kg/m^3}$、$\mathrm{m/s^2}$\\
$d$、$b_c$ & 0.04、0.10 & 坎高、槽宽，m\\
$C_w$ & 1.6 & 交换关系系数，$\mathrm{m^{1/2}/s}$\\
$C_d$、$e$ & 0.62、0.005 & 出流系数、唇口高差（m）\\
原排水口位置 & 1.2、4.8 & m\\
原口有效面积 & $6\times10^{-4}$ & 每口，$\mathrm{m^2}$\\\bottomrule
\end{tabular}
\end{center}

左口位置可取 $0.6、1.2、1.8、2.4\,\mathrm m$，右口可取 $3.6、4.2、4.8、5.4\,\mathrm m$，均自左端量起；各口有效面积可在 $[6,12]\times10^{-4}\,\mathrm{m^2}$ 内连续选取。每改变一个排水口位置增加费用450元，位置不变不收费。扩口费用为
\[
3\times10^6\left[(A_1-6\times10^{-4})+(A_2-6\times10^{-4})\right]\ \text{元},
\]
其中面积代入以 $\mathrm{m^2}$ 表示的数值。可保留原梁，或在全长均匀加固；加固增加自重在本题中忽略。总费用为加固、移口与扩口费用之和。

\begin{center}
\begin{tabular}{crrr}
\toprule
档位 & $EI/EI_0$ & 正负弯矩绝对值限值（N\,m） & 增费（元）\\\midrule
B & 1.00 & 12000 & 0\\
S & 1.25 & 14000 & 3000\\
T & 1.50 & 16000 & 6000\\\bottomrule
\end{tabular}
\end{center}

\begin{center}
\begin{tabular}{p{.42\linewidth}p{.52\linewidth}}
\toprule
文件 & 内容\\\midrule
\file{parameters.json}\newline\file{rules.json} & 物理参数、费用及运行评价规则\\
\file{inflows.csv} & 各工况分段恒定的左右进水量，单位升/分钟；起止时间为分钟\\
\file{stiffening_options.csv} & 加固刚度、弯矩限值与费用\\
\file{data/README.md} & 字段、单位、初始条件与顶撑说明\\\bottomrule
\end{tabular}
\end{center}

\noindent\textbf{参考资料}

\small
积水荷载与结构变形的基础研究可参见 Marino，\href{https://ej.aisc.org/index.php/engj/article/view/1258}{Ponding of Two-Way Roof Systems}，AISC Engineering Journal，1966；Colombi，\textit{The ponding problem on flat steel roof grids}，2006，\href{https://doi.org/10.1016/j.jcsr.2005.11.010}{DOI: 10.1016/j.jcsr.2005.11.010}。实际工程案例可查 \href{https://www.cross-safety.org/uk/safety-information/cross-safety-report/roof-collapse-warehouse-structure-under-rainwater-969}{CROSS：Roof collapse of warehouse structure under rainwater ponding}。本题为原创模拟题，未引用这些资料的试验参数、图表或规范条款。
\end{document}
